\scanpage{48}
forme $4$ fois linéaire sur $E\times\overline E\times F\times\overline F$.
\deleted{\unsure{Dans ceci}} \inserted{Ceci montre} qu'on peut considérer dans $E\otimes F$ \deleted{\unsure{la structure préhilbertienne}} une forme sesquilinéaire $(u,v)$ définie par \deleted{\unsure{une formule}} :
\[
(5)\qquad\boxed{(x_1\otimes y_1,x_2\otimes y_2)=(x_1\deleted{\otimes},x_2)(y_1\deleted{\otimes},y_2)}
\]
(car le deuxième membre de (\unsure{3|5}) est bien linéaire en $x_1\in E$, $x_2\in\overline E$, $y_1\in F$, $y_2\in\overline F$).
\deleted{Montrons que} \inserted{Montrons} qu'on a bien là une forme \unsure{définie positive}.
\deleted{cette forme est une forme définie positive, i.e que $(u,u)\geq0$ et $(u,u)=0\Longleftrightarrow u=0$.
\deleted{En effet} Soit $(e_i)_{i\in I}$ une base orthonormale dans $E$, $(f_j)_{j\in J}$ une base orthonormale dans $F$, alors $(e_i\otimes f_j)_{(ij)\in I\times J}$}

\deleted{En effet, si $\unsure{u}\in E\otimes F$, alors $\unsure{u}$ \unsure{est dans}}
\deleted{$\unsure{u}\in E_0\otimes F_0$, où} \inserted{Donnons-nous} $E_0$ et $F_0$ \deleted{sont} deux sous-espaces de dim finie de $E$ et $F$, \deleted{Prenons alors} \inserted{et prenons} une base orthonormale \inserted{$(e_i)$} dans $E_0$, une autre $(f_j)$ dans $F_0$,
\deleted{\unsure{on aura} $u=\sum$}
alors $(e_i\otimes f_j)$ est une base de l'espace vectoriel, et on sait par (\unsure{5}) que
\[
(6)\qquad(e_i\otimes f_j,e_{i'}\otimes f_{j'})=\delta_{(ij),(i'j')}\qquad\text{(indice de Kronecker)}
\]
ce qui prouve que $(u,v)$ induit sur $E_0\otimes F_0$ une forme \uline{définie positive}.
Comme $E\otimes F$ est réunion des espaces $E_0\otimes F_0$, cela prouve que $(u,v)$ est \uline{définie} positive.

\uline{Définition} Soient $E,F$ deux espaces de Hilbert. On appelle produit tensoriel hilbertien de $E$ et $F$, et on note $E\overset{(2)}{\otimes}F$, le complété de l'espace \deleted{$\unsure{E}$} $E\otimes F$ préhilbertien
\mspage{49}
